% Part: normal-modal-logic
% Chapter: completeness
% Section: modalities-ccs

\documentclass[../../../include/open-logic-section]{subfiles}

\begin{document}

\olfileid{nml}{com}{mod}

\olsection{Modalitas lan Himpunan Konsisten Jangkep}

\begin{explain}
Nalika kita mbangun model~$\mModel{M^\Sigma}$ kang himpunan jagade
kawangun saka himpunan $\Sigma$-konsisten jangkep~$\Delta$ ing
sawijining logika modal normal~$\Sigma$, kita uga kudu netepake
relasi aksesibilitas~$R^\Sigma$ ing antarane ``jagad'' kasebut. Kita
ngarepake relasi aksesibilitas (lan pemetaan $V^\Sigma$) ditepakake
saengga $\mSat{M^\Sigma}{!A}[\Delta]$ yen lan mung yen $!A \in
\Delta$. Kepiye carane?

Sawise relasi aksesibilitas ditepakake, definisi bebener ing
sawijining jagad njamin yen \iftag{prvBox}{$\mSat{M^\Sigma}{\Box!A}[\Delta]$
  yen lan mung yen $\mSat{M^\Sigma}{!A}[\Delta']$ kanggo kabeh $\Delta'$
  kang $R^\Sigma\Delta\Delta'$}{$\mSat{M^\Sigma}{\Diamond!A}[\Delta]$
  yen lan mung yen $\mSat{M^\Sigma}{!A}[\Delta']$ kanggo paling ora siji
  $\Delta'$ kang $R^\Sigma\Delta\Delta'$}. Bukti yen
$\mSat{M^\Sigma}{!A}[\Delta]$ yen lan mung yen $!A \in \Delta$
mbutuhake iki uga bener tumrap !!{formula}s kang diawali operator modal,
yaiku \iftag{prvBox}{$\mSat{M^\Sigma}{\Box!A}[\Delta]$ yen lan mung yen
  $\Box !A \in \Delta$}{$\mSat{M^\Sigma}{\Diamond!A}[\Delta]$ yen lan mung yen
  $\Diamond !A \in \Delta$}. Yen sarat iki digandhengake karo definisi
bebener ing sawijining jagad kanggo \iftag{prvBox}{$\Box !A$}{$\Diamond !A$},
kita entuk:
\[
\iftag{prvBox}{%
  \Box!A \in \Delta \text{ yen lan mung yen } !A \in \Delta' \text{
    kanggo kabeh $\Delta'$ kanthi } R^\Sigma\Delta\Delta'}{%
  \Diamond!A \in \Delta \text{ yen lan mung yen } !A \in \Delta' \text{ kanggo
    paling ora siji } \Delta' \text{ kanthi } R^\Sigma\Delta\Delta'}
\]
\iftag{prvBox}{Delengen arah kiwa menyang tengen: yen
  $\Box !A \in \Delta$, mula $!A \in \Delta'$ kanggo sembarang $!A$ lan
  sembarang $\Delta'$ kang $R^\Sigma\Delta\Delta'$. Yen kita netepake
  $R^\Sigma\Delta\Delta'$ yen lan mung yen $!A \in \Delta'$ kanggo kabeh
  $\Box!A \in \Delta$, mula sarat iku kawujud. Sarat ing sisih tengen
  saka ``yen lan mung yen'' bisa ditulis luwih ringkes minangka
  $\Setabs{!A}{\Box!A \in \Delta} \subseteq \Delta'$.}{Delengen arah
  tengen menyang kiwa: yen $!A \in \Delta'$, mula $\Diamond!A \in \Delta$,
  kanggo sembarang $!A$ lan $\Delta'$ kang $R^\Sigma\Delta\Delta'$.
  Yen kita netepake $R^\Sigma\Delta\Delta'$ yen lan mung yen
  $\Diamond!A \in \Delta$ kanggo kabeh $!A \in \Delta'$, sarat iku
  kawujud. Sarat ing sisih tengen saka ``yen lan mung yen'' bisa ditulis
  luwih ringkes minangka
  $\Setabs{\Diamond!A}{!A \in \Delta'} \subseteq \Delta$.}

Dadi pitakone: apa definisi $R^\Sigma$ iki pancen njamin
yen \iftag{prvBox}{$\Box !A \in \Delta$ yen lan mung yen
  $\mSat{M^\Sigma}{\Box !A}[\Delta]$? \iftag{prvDiamond}{Apa uga
    njamin yen }{}}{}\iftag{prvDiamond}{$\Diamond !A \in \Delta$
  yen lan mung yen $\mSat{M^\Sigma}{\Diamond !A}[\Delta]$?}{}
Asil-asil sabanjure bakal netepake iki.
\end{explain}

\begin{defn}
  Yen $\Gamma$ himpunan !!{formula}s, tetepa
  \begin{align*}
    \Box\Gamma & = \Setabs{\Box !B}{!B \in \Gamma}\\
    \Diamond\Gamma & = \Setabs{\Diamond !B}{!B \in \Gamma}\\
    \intertext{lan}
    \Box^{-1}\Gamma & = \Setabs{!B}{\Box !B \in \Gamma}\\
    \Diamond^{-1}\Gamma & = \Setabs{!B}{\Diamond !B \in \Gamma}\\
  \end{align*}
\end{defn}

Kanthi tembung liya, $\Box\Gamma$ iku $\Gamma$ kanthi
$\Box$ ing ngarepe saben !!{formula} ing~$\Gamma$;
$\Box^{-1}\Gamma$ iku himpunan kabeh !!{formula}s kang diawali $\Box$
ing~$\Gamma$ sawisé $\Box$ wiwitane dicopot. Definisi iki
ora wigati banget kanthi mandiri, nanging bakal nggampangake notasi.

Elinga yen $\Box\Box^{-1}\Gamma \subseteq \Gamma$:
\[
\Box\Box^{-1}\Gamma = \Setabs{\Box!B}{\Box!B \in \Gamma}
\]
yaiku himpunan kabeh !!{formula}s ing~$\Gamma$ kang diawali~$\Box$.

\begin{lem}\ollabel{lem:box1}
  Yen $\Gamma \Proves[\Sigma] !A$, mula $\Box\Gamma \Proves[\Sigma] \Box
  !A$.
\end{lem}

\begin{proof}
  Yen $\Gamma \Proves[\Sigma] !A$, mula ana $!B_1$, \dots, $!B_k
  \in \Gamma$ saengga $\Sigma \Proves !B_1 \lif (!B_2 \lif \cdots
  (!B_k \lif !A)\cdots)$. Amarga $\Sigma$ normal, miturut aturan \RK{},
  $\Sigma \Proves \Box!B_1 \lif (\Box!B_2 \lif \cdots (\Box!B_k \lif
  \Box!A)\cdots)$; lan cetha yen $\Box!B_1$, \dots, $\Box!B_k \in
  \Box\Gamma$. Mula, miturut definisi, $\Box\Gamma \Proves[\Sigma] \Box
  !A$.
\end{proof}

\begin{lem}\ollabel{lem:box2}
  Yen $\Box^{-1} \Gamma \Proves[\Sigma] !A$, mula $\Gamma
  \Proves[\Sigma] \Box!A$.
\end{lem}

\begin{proof}
  Upamane $\Box^{-1}\Gamma \Proves[\Sigma] !A$; banjur miturut
  \olref{lem:box1}, $\Box\Box^{-1}\Gamma \Proves[\Sigma] \Box!A$.
  Nanging amarga $\Box\Box^{-1}\Gamma \subseteq \Gamma$, miturut
  monotonisitas uga $\Gamma \Proves[\Sigma] \Box!A$.
\end{proof}

\iftag{prvBox}{%
% Box is primitive

\begin{prop}\ollabel{prop:box}
  Yen $\Gamma$ $\Sigma$-konsisten jangkep, mula $\Box !A \in
  \Gamma$ yen lan mung yen kanggo saben himpunan $\Sigma$-konsisten
  jangkep $\Delta$ kanthi $\Box^{-1} \Gamma \subseteq \Delta$, kita entuk
  $!A \in \Delta$.
\end{prop}

\begin{proof}
  Upamane $\Gamma$ $\Sigma$-konsisten jangkep. Arah ``mung yen''
  gampang: upamane $\Box !A \in \Gamma$ lan
  $\Box^{-1}\Gamma \subseteq \Delta$. Amarga $\Box!A \in \Gamma$, $!A
  \in \Box^{-1}\Gamma \subseteq \Delta$, mula $!A \in \Delta$.

Kanggo arah ``yen'', kita mbuktekake kontraposisine: upamane
  $\Box!A \notin \Gamma$. Amarga $\Gamma$ $\Sigma$-konsisten
  jangkep, himpunan iku katutup sacara deduktif, mula $\Gamma
  \Proves/[\Sigma] \Box !A$. Miturut \olref{lem:box2},
  $\Box^{-1}\Gamma \Proves/[\Sigma] !A$. Miturut
  \olref[prf][con]{prop:consistencyfacts}\olref[prf][con]{prop:consistencyfacts-b},
  $\Box^{-1}\Gamma \cup \{ \lnot!A \}$ iku $\Sigma$-konsisten. Miturut
  Lemma Lindenbaum, ana himpunan $\Sigma$-konsisten jangkep~$\Delta$
  saengga $\Box^{-1}\Gamma \cup \{ \lnot!A \} \subseteq
  \Delta$. Amarga konsistensi, $!A \notin \Delta$.
\end{proof}
}{%
% Box is not primitive, only Diamond is

\begin{prop}\ollabel{prop:diamond}
  Yen $\Gamma$ $\Sigma$-konsisten jangkep, mula $\Diamond !A \in
  \Gamma$ yen lan mung yen kanggo sawenehing himpunan $\Sigma$-konsisten
  jangkep $\Delta$ kanthi $\Diamond\Delta \subseteq
  \Gamma$, kita entuk $!A \in \Delta$.
\end{prop}

\begin{proof}
  Upamane $\Gamma$ $\Sigma$-konsisten jangkep. Arah tengen menyang kiwa
  gampang: pilih $\Delta$ $\Sigma$-konsisten jangkep kanthi
  $\Diamond\Delta \subseteq \Gamma$ lan $!A \in \Delta$. Banjur
  $\Diamond!A \in \Diamond\Delta \subseteq \Gamma$, yaiku $\Diamond!A
  \in \Gamma$.

Kanggo arah kiwa menyang tengen, upamane $\Diamond !A \in \Gamma$.
  Kita kudu nuduhake anane himpunan $\Sigma$-konsisten jangkep $\Delta$
  kanthi $!A \in \Delta$ lan $\Diamond\Delta \subseteq \Gamma$. Amarga
  $\Diamond !A \in \Gamma$ lan $\Gamma$ $\Sigma$-konsisten,
  $\Box\lnot !A \notin \Gamma$. Amarga $\Gamma$ $\Sigma$-konsisten
  jangkep, himpunan iki katutup sacara deduktif, mula $\Gamma
  \Proves/[\Sigma] \Box\lnot !A$. Miturut \olref{lem:box2},
  $\Box^{-1}\Gamma \Proves/[\Sigma] \lnot !A$. Miturut
  \olref[prf][con]{prop:consistencyfacts}\olref[prf][con]{prop:consistencyfacts-b},
  $\Box^{-1}\Gamma \cup \{ !A \}$ iku $\Sigma$-konsisten. Miturut
  Lemma Lindenbaum, ana himpunan $\Sigma$-konsisten jangkep~$\Delta$
  saengga $\Box^{-1}\Gamma \cup \{ !A \} \subseteq
  \Delta$. Cetha yen $!A \in \Delta$. Kanggo nuduhake $\Diamond\Delta
  \subseteq \Gamma$, upamane ora mangkono: kanggo sawenehing $!B \in \Delta$,
  $\Diamond !B \notin \Gamma$. Amarga $\Gamma$ $\Sigma$-konsisten
  jangkep, $\Box\lnot !B \in \Gamma$. Nanging uga $\lnot
  !B \in \Box^{-1}\Gamma \subseteq \Delta$, saengga $\Delta$
  ora konsisten.
\end{proof}
}

\begin{lem}\ollabel{lem:box-iff-diamond}
  Upamane $\Gamma$ lan $\Delta$ $\Sigma$-konsisten jangkep.
  Banjur $\Box^{-1}\Gamma \subseteq \Delta$ yen lan mung yen
  $\Diamond\Delta \subseteq \Gamma$.
\end{lem}

\begin{proof}
  Arah ``mung yen'': upamane $\Box^{-1}\Gamma \subseteq \Delta$ lan
  $\Diamond!A \in \Diamond\Delta$ (yaiku $!A \in \Delta$). Kanggo
  nuduhake $\Diamond!A \in \Gamma$, cukup nuduhake yen
  $\Box\lnot!A \notin \Gamma$, amarga miturut maksimalitas,
  $\lnot\Box\lnot !A \in \Gamma$. Saiki, yen $\Box\lnot!A \in \Gamma$,
  miturut hipotesis $\lnot!A \in \Delta$, kang mbantah konsistensi
  $\Delta$ (amarga $!A \in \Delta$). Mula $\Box\lnot!A \notin \Gamma$,
  kaya kang dibutuhake.

Arah ``yen'': upamane $\Diamond\Delta
  \subseteq \Gamma$. Kita nggunakake kontraposisi: upamane $!A
  \notin \Delta$ kanggo nuduhake $\Box!A \notin \Gamma$. Yen
  $!A \notin \Delta$, miturut maksimalitas $\lnot!A \in \Delta$,
  mula miturut hipotesis $\Diamond\lnot!A \in \Gamma$. Nanging ing
  logika modal normal, $\Diamond\lnot!A$ ekuivalen karo
  $\lnot\Box !A$, lan yen kang pungkasan iki ana ing $\Gamma$, miturut
  konsistensi $\Box!A \notin\Gamma$, kaya kang dibutuhake.
\end{proof}

\iftag{prvBox}{%
  \iftag{prvDiamond}{%
% Box and Diamond are both primitive; prove
% prop:diamond using lem:box-iff-diamond

\begin{prop}\ollabel{prop:diamond}
  Yen $\Gamma$ $\Sigma$-konsisten jangkep, mula $\Diamond !A \in
  \Gamma$ yen lan mung yen kanggo sawenehing himpunan $\Sigma$-konsisten
  jangkep $\Delta$ kanthi $\Diamond\Delta \subseteq
  \Gamma$, kita entuk $!A \in \Delta$.
\end{prop}

\begin{proof}
Upamane $\Gamma$ $\Sigma$-konsisten jangkep. $\Diamond!A \in
\Gamma$ yen lan mung yen $\lnot\Box\lnot !A \in \Gamma$ miturut \Dual{}
lan katutup sacara deduktif. $\lnot\Box\lnot !A \in \Gamma$ yen lan
mung yen $\Box\lnot !A \notin \Gamma$ miturut
\olref[ccs]{prop:ccs-properties}\olref[ccs]{prop:ccs-lnot} amarga
$\Gamma$ $\Sigma$-konsisten jangkep. Miturut \olref{prop:box},
$\Box\lnot !A \notin \Gamma$ yen lan mung yen ana himpunan
$\Sigma$-konsisten jangkep $\Delta$ kanthi $\Box^{-1}\Gamma \subseteq \Delta$
lan $\lnot!A \notin \Delta$. Saiki delengen sembarang~$\Delta$ kaya
mangkono. Miturut \olref{lem:box-iff-diamond},
$\Box^{-1}\Gamma \subseteq \Delta$ yen lan mung yen
$\Diamond\Delta \subseteq \Gamma$. Uga, $\lnot !A \notin \Delta$ yen lan
mung yen $!A \in \Delta$ miturut
\olref[ccs]{prop:ccs-properties}\olref[ccs]{prop:ccs-lnot}. Dadi
$\Diamond!A \in \Gamma$ yen lan mung yen ana himpunan
$\Sigma$-konsisten jangkep $\Delta$ kanthi $\Diamond\Delta \subseteq \Gamma$
lan $!A \in \Delta$.
\end{proof}

}{}}{}

\begin{prob}
  Tuduhna yen $\Gamma$ $\Sigma$-konsisten jangkep, mula
  \iftag{prvBox}{$\Diamond !A \in \Gamma$ yen lan mung yen ana
    himpunan $\Sigma$-konsisten jangkep $\Delta$ kanthi $\Box^{-1}\Gamma
    \subseteq \Delta$ lan $!A \in \Delta$.}{$\Box !A \in
    \Gamma$ yen lan mung yen kanggo saben himpunan $\Sigma$-konsisten
    jangkep $\Delta$ kanthi $\Box^{-1} \Gamma \subseteq \Delta$, kita entuk
    $!A \in \Delta$.} \emph{Lakonana tanpa nggunakake
    \olref[nml][com][mod]{lem:box-iff-diamond}}.
\end{prob}

\end{document}
